\scanpage{16}
\mspage{-0.18-}

one $x$. To illustrate this, we give an example and a lemma which suggest opposite conclusions (the example shattered over-optimistic conjectures that the author had entertained!).

\uline{Example} \uline{0.4}: Of a semi-simple group $G$ over $\handwritten{\mathbb C}$, the complex numbers, acting on a non-singular quasi-affine scheme of finite type over $\handwritten{\mathbb C}$, such that:

i) all stabilizers are reduced to $e$ itself, i.e. the action is set-theoretically \tikz[baseline=(word.base)]{\node[ellipse,draw,inner sep=2pt] (word) {free,};}
% Pencil oval around the typed word and comma.

ii) a geometric quotient exists, which is also a scheme of finite type over $\handwritten{\mathbb C}$; hence (cf. \S4) the action is separated,

\margin{\handwritten{\unsure{$\sim$!|!}}}
iii) the action is not proper, and, in particular, \uline{not\unsure{:|;|i:} free} (algebro-geometrically).

\uline{Proof}: Take $G=\mathrm{SL}(2)$. Let $V_n$ stand for the $(n+1)$ -dimensional affine space whose closed points are homogeneous forms in $X$ and $Y$ of degree $n$ over $\handwritten{\mathbb C}$ : then $\mathrm{SL}(2)$ acts on $V_n$ by means of substitutions in $X$ and $Y$. Define a reduced subscheme $X\handwritten{\subset}V_1\times V_4$ by means of
\[
\begin{aligned}
(F_1,F_4)\in X_{\handwritten{\mathbb C}}\handwritten{\Longleftrightarrow} &\quad\text{(a) }F_1\ne 0\\
 &\quad\text{(b) }F_4\ \text{ is the square of a homogeneous}\\
 &\qquad\quad\text{quadratic form of discriminant 1.}
\end{aligned}
\]
One checks that $X$ is non-singular, invariant under $\mathrm{SL}(2)$, and that no non-trivial $\handwritten{\alpha}\in\mathrm{SL}(2)_{\handwritten{\mathbb C}}$ leaves fixed any point of $X$. Define $\unsure{\phi|p}:X\handwritten{\to\mathbb A}^{1}$ by
% archive sheet number: 15

